\documentclass[11pt,a4paper]{article}

\usepackage{../shared/cathedral-preamble}
\usepackage{setspace}
\onehalfspacing

\title{%
  The Dark Mirror:\\
  Dual-Use Risks of Number-Theoretic Engineering\\
  from the Cathedral Project%
}
\author{
  Jason Robert Gochanour
}
\date{April 20, 2026}

\begin{document}

\maketitle

\begin{abstract}
The companion papers \emph{What Could Be Built} and \emph{The
Prime Number Grid} propose engineering applications of the
mathematical structures discovered during the Cathedral project.
Responsible disclosure requires examining the same technologies
from an adversarial perspective. This paper identifies five
dual-use risks: (1)~precision harmonic injection attacks on power
grids, (2)~resonance-targeted industrial sabotage via arithmetic
vibration analysis, (3)~adversarial use of verified instability
certificates for cascade failure engineering, (4)~broadband stealth
and jamming from prime-spaced antenna arrays, and (5)~covert RF
surveillance via arithmetic energy harvesting. For each risk, we
assess the severity, the novelty (whether the attack vector is
genuinely new), the current state of defenses, and recommended
mitigations. We conclude that the mathematical structures in the
Cathedral do not create fundamentally new attack capabilities, but
could make existing attack vectors more \emph{precise} and harder
to detect. We recommend specific defensive measures and argue that
publishing this analysis openly is the responsible path: attackers
already have access to these mathematical tools; defenders need to
understand them too. This paper exists because the authors believe
that \emph{thinking about misuse before deployment} is a moral
obligation of anyone who builds powerful tools.
\end{abstract}

\tableofcontents

% ================================================================
\section{Why This Paper Exists}
% ================================================================

Every powerful tool has a dual use. The same chemistry that produces
fertilizer produces explosives. The same physics that enables
nuclear medicine enables nuclear weapons. The same number theory
that could improve power grid monitoring could, in principle, be
turned against the grid.

The companion papers propose six positive energy applications. This
paper examines what happens when each is reflected through an
adversarial lens. We do not provide implementation details for any
attack---we describe the \emph{category} of risk and the
\emph{structural} vulnerability, not the recipe.

\subsection{Responsible Disclosure Rationale}

We publish this analysis openly for three reasons:

\begin{enumerate}
  \item \textbf{Attackers already know the math.} The M\"obius
    function, Fourier analysis, spectral theory, and variational
    methods are textbook material. Anyone with a graduate education
    in mathematics or engineering already has access to these tools.
    Suppressing this paper would not deny attackers any capability.

  \item \textbf{Defenders need to understand the threats.} Grid
    operators, infrastructure security teams, and turbine
    manufacturers benefit from understanding how number-theoretic
    tools could be used against their systems. Silence helps
    attackers more than defenders.

  \item \textbf{Moral obligation.} If we propose technologies, we
    must also assess their misuse potential. Publishing only the
    positive applications would be intellectually dishonest and
    ethically incomplete.
\end{enumerate}

% ================================================================
\section{Risk 1: Precision Harmonic Injection Attacks on Power Grids}
\label{sec:harmonics}
% ================================================================

\subsection{The Positive Technology}

The M\"obius Harmonic Analyzer (MHA) decomposes grid harmonics into
independent sources, enabling targeted active filtering.

\subsection{The Dark Mirror}

The same M\"obius decomposition that identifies harmonic sources
can be used to \emph{synthesize} optimal harmonic injections. An
attacker with a sufficiently powerful inverter connected to the
grid could:

\begin{enumerate}
  \item \textbf{Targeted transformer saturation}: Inject
    even-order harmonics (2nd, 4th) that cause DC offset in the
    transformer core flux, leading to half-cycle saturation.
    The M\"obius decomposition identifies exactly which harmonic
    frequencies will bypass existing passive filters.

  \item \textbf{Capacitor bank resonance amplification}: Identify
    the resonant frequency of grid capacitor banks
    ($f_r = 1/(2\pi\sqrt{LC})$) and inject power at that frequency.
    If $f_r$ is near an integer harmonic, the arithmetic structure
    of the injection maximizes energy transfer.

  \item \textbf{Protective relay confusion}: Craft a waveform whose
    zero-crossing pattern mimics a fault condition, causing
    protective relays to trip unnecessarily. The M\"obius-optimal
    harmonic mixture minimizes the injected power needed to trigger
    a trip.

  \item \textbf{Stealth}: By using M\"obius inversion, the attacker
    can design an injection that appears as ``natural'' harmonic
    distortion (consistent with normal nonlinear load behavior)
    rather than an obvious attack signal.
\end{enumerate}

\subsection{Severity Assessment}

\begin{center}
\begin{tabular}{@{}ll@{}}
\toprule
\textbf{Factor} & \textbf{Assessment} \\
\midrule
Novelty & Low --- harmonic injection is a known attack vector \\
Precision gain & Medium --- M\"obius optimizes the injection \\
Stealth gain & \textbf{High} --- attack looks like natural distortion \\
Physical access required & Grid connection (industrial inverter) \\
Detection difficulty & High if designed using M\"obius structure \\
Potential impact & Localized transformer/capacitor damage \\
\bottomrule
\end{tabular}
\end{center}

\subsection{Defenses}

\begin{enumerate}
  \item \textbf{Deploy the MHA defensively}: The same tool that
    enables the attack also detects it. An MHA-equipped grid monitor
    can identify ``unnatural'' harmonic source patterns that don't
    match known load profiles.
  \item \textbf{Harmonic injection limits}: IEEE~519 already sets
    limits on harmonic current injection. Enforce these at the
    point of common coupling with real-time monitoring.
  \item \textbf{Anomaly detection}: Train ML models on the
    M\"obius-decomposed harmonic spectrum to detect injections that
    deviate from historical patterns.
\end{enumerate}

% ================================================================
\section{Risk 2: Resonance-Targeted Industrial Sabotage}
\label{sec:vibration}
% ================================================================

\subsection{The Positive Technology}

The Arithmetic Vibration Monitor (AVM) uses compressed sensing with
an arithmetic dictionary to detect incipient turbine faults.

\subsection{The Dark Mirror}

If arithmetic compressed sensing can detect subtle vibration
signatures, it can also \emph{design} them. An attacker with
physical access to a turbine (or the ability to modulate its load)
could:

\begin{enumerate}
  \item \textbf{Resonance excitation}: Identify the natural
    frequencies of turbine blades via the Gram matrix eigenvalue
    structure, then apply periodic load variations at a blade
    natural frequency to induce resonant fatigue failure.

  \item \textbf{Sub-harmonic injection}: Excite blade vibration at
    a frequency that is $1/k$ of a natural frequency ($k$ prime),
    so the vibration builds gradually through nonlinear coupling
    rather than direct resonance. This is harder to detect because
    the excitation frequency is not near any monitored harmonic.

  \item \textbf{Monitoring evasion}: Design the excitation signal
    to be orthogonal to the AVM's sensing dictionary, making it
    invisible to the monitoring system while still delivering
    destructive energy to the blade.
\end{enumerate}

\subsection{Severity Assessment}

\begin{center}
\begin{tabular}{@{}ll@{}}
\toprule
\textbf{Factor} & \textbf{Assessment} \\
\midrule
Novelty & Medium --- resonance attacks exist, arithmetic targeting is new \\
Precision gain & \textbf{High} --- spectral analysis identifies exact modes \\
Stealth gain & \textbf{High} --- sub-harmonic excitation evades monitoring \\
Physical access required & Load modulation or physical access to nacelle \\
Detection difficulty & Very high for sub-harmonic approach \\
Potential impact & Blade failure, turbine destruction, potential casualties \\
\bottomrule
\end{tabular}
\end{center}

\subsection{Defenses}

\begin{enumerate}
  \item \textbf{Broadband monitoring}: Monitor vibration across a
    wider frequency range than just the expected harmonics.
    Specifically monitor sub-harmonics at $f_{\text{nat}}/p$ for
    prime $p$.
  \item \textbf{Load pattern anomaly detection}: Flag unusual load
    variation patterns, especially periodic patterns at non-integer
    frequencies.
  \item \textbf{Physical security}: Restrict physical access to
    turbine nacelles and control systems. This is the most effective
    defense against all vibration-based attacks.
\end{enumerate}

% ================================================================
\section{Risk 3: Verified Instability Certificates}
\label{sec:instability}
% ================================================================

\subsection{The Positive Technology}

The Cathedral's framework produces machine-verified stability
certificates for grid configurations: formal proofs that a
Lyapunov function is positive definite.

\subsection{The Dark Mirror}

The same framework can produce machine-verified \emph{instability}
certificates: formal proofs that a specific perturbation sequence
causes cascade failure.

\begin{enumerate}
  \item \textbf{Optimal cascade sequences}: Given a grid model,
    find the minimum set of generator trips that causes voltage
    collapse. The Lean~4 verification ensures the attack plan is
    mathematically correct --- no ``wasted'' disruptions.

  \item \textbf{Stability boundary mapping}: Identify the exact
    operating conditions where the grid transitions from stable to
    unstable. An attacker who can push the grid to this boundary
    (e.g., by coordinating large loads) can trigger instability with
    minimal visible action.

  \item \textbf{Targeted N-k analysis}: Current grid security
    standards require N-1 or N-2 contingency analysis (survive the
    loss of 1 or 2 components). A verified instability certificate
    could identify the specific N-3 or N-4 scenarios that current
    analysis misses.
\end{enumerate}

\subsection{Severity Assessment}

\begin{center}
\begin{tabular}{@{}ll@{}}
\toprule
\textbf{Factor} & \textbf{Assessment} \\
\midrule
Novelty & \textbf{High} --- verified attack plans are new \\
Precision gain & \textbf{Very High} --- mathematically optimal attack \\
Stealth gain & Medium --- the attack itself is visible \\
Access required & Grid topology and parameters (often semi-public) \\
Detection difficulty & Low (the cascade is observable) \\
Potential impact & \textbf{Regional blackout, economic damage} \\
\bottomrule
\end{tabular}
\end{center}

This is the most concerning risk in this paper. The novelty is not
in the attack (deliberate grid disruption is well-studied) but in
the \textbf{verification}: a machine-checked proof that the attack
\emph{will} succeed, removing uncertainty from the attacker's
planning.

\subsection{Defenses}

\begin{enumerate}
  \item \textbf{Use the same tool defensively}: If the attacker can
    find verified instabilities, the grid operator should find them
    first. Deploy the same Lean~4 framework to systematically
    identify and remediate all N-3 and N-4 vulnerabilities.
  \item \textbf{Protect grid models}: Limit public access to
    detailed grid topology and impedance data. Current CEII
    (Critical Energy Infrastructure Information) rules should be
    reviewed and strengthened.
  \item \textbf{Dynamic reconfiguration}: Design grids that can
    rapidly change topology in response to detected instability,
    invalidating pre-computed attack plans.
  \item \textbf{Adversarial stability margins}: Design with larger
    stability margins specifically in regions where N-3/N-4
    analysis reveals verified instabilities.
\end{enumerate}

% ================================================================
\section{Risk 4: Prime-Spaced Stealth and Jamming}
\label{sec:stealth}
% ================================================================

\subsection{The Positive Technology}

Prime-spaced sensor arrays with Selberg beamforming weights offer
broadband performance with predictable sidelobe control.

\subsection{The Dark Mirror}

The same array design, used in reverse, can produce electromagnetic
emissions that are:

\begin{enumerate}
  \item \textbf{Broadband noise with minimal power}: A prime-spaced
    jammer distributes its energy across frequencies without
    periodic structure, making it harder to filter. Conventional
    interference cancellation relies on identifying the jammer's
    spectral structure; a prime-spaced jammer has no periodic
    structure to identify.

  \item \textbf{Difficult to geolocate}: The aperiodic radiation
    pattern of a prime-spaced array produces a complex spatial
    pattern that conventional direction-finding algorithms
    (designed for uniform or periodic arrays) may fail to resolve.

  \item \textbf{Radar-evasive scattering}: A target coated in
    prime-spaced resonant elements would scatter radar returns
    across frequencies without the signature Bragg peaks of a
    periodic structure, reducing radar cross-section.
\end{enumerate}

\subsection{Severity Assessment}

\begin{center}
\begin{tabular}{@{}ll@{}}
\toprule
\textbf{Factor} & \textbf{Assessment} \\
\midrule
Novelty & Low --- aperiodic jamming/stealth are well-studied \\
Precision gain & Low --- random spacing works comparably \\
Stealth gain & Low --- prime spacing has marginal advantage \\
Access required & RF hardware (commercially available) \\
Detection difficulty & Medium \\
Potential impact & Communications disruption \\
\bottomrule
\end{tabular}
\end{center}

\textbf{Honest assessment}: This is the weakest risk in the paper.
Prime-spaced arrays are marginally better than random arrays for
jamming/stealth purposes, and the advantage does not justify
concern beyond existing aperiodic jamming threats.

\subsection{Defenses}

Existing anti-jamming techniques (spread spectrum, frequency
hopping, adaptive beamforming) are effective against aperiodic
jammers regardless of whether the aperiodicity comes from random
or prime spacing.

% ================================================================
\section{Risk 5: Covert Broadband RF Surveillance}
\label{sec:surveillance}
% ================================================================

\subsection{The Positive Technology}

Prime-indexed rectennas harvest ambient RF energy across a wide
bandwidth, enabling battery-free IoT sensors.

\subsection{The Dark Mirror}

A device that harvests broadband RF also \emph{receives} broadband
RF. A prime-indexed rectenna array disguised as an energy harvester
could:

\begin{enumerate}
  \item \textbf{Passive SIGINT}: Simultaneously monitor WiFi
    (2.4/5 GHz), cellular (700 MHz--3 GHz), Bluetooth (2.4 GHz),
    and IoT (900 MHz) without active emissions, making it
    undetectable by RF sweeps that look for transmitters.

  \item \textbf{Self-powered}: The same harvested energy that
    powers the device also carries the intelligence. No battery
    means no maintenance visits, no supply chain to trace, and
    indefinite operational life.

  \item \textbf{Plausible deniability}: If discovered, the device
    appears to be a legitimate energy harvester (IoT sensor, smart
    building component) rather than a surveillance device.
\end{enumerate}

\subsection{Severity Assessment}

\begin{center}
\begin{tabular}{@{}ll@{}}
\toprule
\textbf{Factor} & \textbf{Assessment} \\
\midrule
Novelty & Medium --- passive RF surveillance exists, self-powered is new \\
Precision gain & Low \\
Stealth gain & \textbf{High} --- disguised as legitimate IoT hardware \\
Access required & Physical placement in target facility \\
Detection difficulty & \textbf{Very high} --- no RF emissions to detect \\
Potential impact & Espionage, corporate/government intelligence \\
\bottomrule
\end{tabular}
\end{center}

\subsection{Defenses}

\begin{enumerate}
  \item \textbf{Physical inspection}: Regular RF environment audits
    that include physical inspection of all connected devices.
  \item \textbf{Faraday shielding}: Sensitive facilities should
    control RF ingress/egress regardless of harvesting technology.
  \item \textbf{Supply chain security}: Verify the provenance of
    all IoT/energy harvesting devices deployed in sensitive
    environments.
\end{enumerate}

% ================================================================
\section{Risk Assessment Summary}
% ================================================================

\begin{center}
\small
\begin{tabular}{@{}lcccl@{}}
\toprule
\textbf{Risk} & \textbf{Novelty} & \textbf{Severity} &
\textbf{Feasibility} & \textbf{Primary Defense} \\
\midrule
Grid harmonic injection & Low & Medium & High & MHA defensive monitoring \\
Vibration sabotage & Medium & High & Medium & Broadband monitoring \\
Verified instability & \textbf{High} & \textbf{Very High} &
  Medium & Defensive verification \\
Stealth/jamming & Low & Low & High & Existing countermeasures \\
RF surveillance & Medium & Medium & Medium & Physical/supply chain \\
\bottomrule
\end{tabular}
\end{center}

\subsection{The Headline Risk}

The single most concerning capability is \textbf{Risk~3: Verified
Instability Certificates}. The ability to produce a
\emph{machine-checked proof} that a specific attack sequence will
cause a grid cascade failure is qualitatively different from
simulation-based attack planning. A simulation can be wrong; a
verified certificate cannot (within its axioms).

The mitigation is symmetric: grid operators should use the same
framework to find these vulnerabilities first. The race is between
defensive and offensive verification, and the defender has the
advantage of knowing their own system.

% ================================================================
\section{The Broader Dual-Use Question}
% ================================================================

None of the risks identified in this paper are created by the
Cathedral. The mathematical tools (Fourier analysis, spectral
theory, M\"obius function, variational methods) have been publicly
available for decades to centuries. What the Cathedral contributes
is:

\begin{enumerate}
  \item \textbf{Organization}: The Cathedral assembles these tools
    into a coherent framework. An attacker would otherwise need to
    independently recognize the connections between number theory
    and engineering.
  \item \textbf{Verification}: The machine-verified proof
    methodology could, in principle, be applied to verify attack
    plans rather than proofs.
  \item \textbf{Precision}: The number-theoretic structure (primes,
    M\"obius function, Selberg weights) provides optimal or
    near-optimal parameters for certain engineering tasks,
    including adversarial ones.
\end{enumerate}

The dual-use question is not whether to publish (the math is
already public) but whether to \emph{emphasize the connections}
between number theory and infrastructure engineering. We believe
the answer is yes, because:

\begin{quote}
\emph{The asymmetry in dual-use technology favors the defender.}
\end{quote}

Defenders (grid operators, turbine manufacturers, infrastructure
security teams) have domain knowledge, system access, and
institutional resources. Attackers must operate covertly, with
limited access and imperfect models. Providing defenders with a
new mathematical framework strengthens them more than it
strengthens attackers.

% ================================================================
\section{Recommendations}
% ================================================================

\begin{enumerate}
  \item \textbf{Deploy the MHA defensively before it is used
    offensively.} The M\"obius Harmonic Analyzer is the most
    immediately actionable tool in both the positive and negative
    applications. Grid operators should adopt it for monitoring
    before attackers adopt it for injection design.

  \item \textbf{Extend N-k contingency analysis.} Use the Cathedral's
    verified stability framework to systematically identify N-3 and
    N-4 grid vulnerabilities. Remediate them before an adversary
    discovers them independently.

  \item \textbf{Add sub-harmonic channels to vibration monitors.}
    Current turbine monitoring focuses on integer harmonics of the
    rotation frequency. Adding monitoring at $f/p$ for small primes
    $p$ would detect the sub-harmonic excitation attack described
    in Risk~2.

  \item \textbf{Classify verified instability certificates.}
    Machine-checked proofs of grid instability should be treated
    as sensitive information (CEII) and protected accordingly.

  \item \textbf{Publish this paper.} Suppressing dual-use analysis
    protects no one. The mathematical tools are public. The
    infrastructure is public. The only thing that can be private is
    the defender's understanding of their own vulnerabilities---and
    that understanding is strengthened, not weakened, by open
    analysis.
\end{enumerate}

% ================================================================
\section{Conclusion}
% ================================================================

The Cathedral's mathematics is neutral. The Parseval Bridge connects
time and frequency domains. The M\"obius function decomposes
multiplicative structure. The Rayleigh quotient bounds eigenvalues.
These tools do not care whether they are used to protect a power
grid or to attack one.

What matters is who learns to use them first.

This paper exists to ensure that the people who protect critical
infrastructure understand the same mathematics that could threaten
it. We believe that open, honest analysis of dual-use risks---even
uncomfortable analysis---is a more effective defense than silence.

The Cathedral was built to illuminate. This paper points the light
at the shadows.

\begin{quote}
\emph{A tool that can diagnose can also deceive.\\
A proof that certifies stability can also certify its absence.\\
The mathematics does not choose sides.\\
We must.}
\end{quote}

\vspace{0.5cm}

\begin{thebibliography}{99}

\bibitem{ieee519}
IEEE, \emph{IEEE Std 519-2022: Recommended Practice for Harmonic
Control}, 2022.

\bibitem{nerc}
NERC, \emph{Grid Security Guidelines: Critical Infrastructure
Protection Standards CIP-002 through CIP-014}, 2023.

\bibitem{stuxnet}
R.~Langner, \emph{Stuxnet: Dissecting a Cyberwarfare Weapon},
IEEE Security \& Privacy, 2011.

\bibitem{ceii}
FERC, \emph{Critical Energy Infrastructure Information (CEII)
Regulations}, 18 CFR \S 388.113.

\bibitem{cathedral-futures}
J.R.~Gochanour, \emph{What Could Be Built: Seven Engineering
Frontiers}, 2026.

\bibitem{cathedral-energy}
J.R.~Gochanour, \emph{The Prime Number Grid: Energy Systems
Applications}, 2026.

\end{thebibliography}

\end{document}
